% =====================================================================================
% REVISIÓN 1: solo los «Temas tentativos del marco teórico».
% REVISIÓN 2: el Marco teórico completo (cuando lo desarrolles, sustituye o convierte los temas tentativos en subsecciones).
% Escribe tu texto FUERA de los recuadros \begin{guia} ... \end{guia}.
% =====================================================================================

\section{Marco teórico} \label{sec:marco_teorico} % audit:req=marco_teorico

A continuación se presentan los temas tentativos que se desarrollarán en este capítulo. Se organizan por concepto y no por herramienta: las plataformas concretas (Power Platform, Power BI, SharePoint, SAP) se describen en el capítulo de desarrollo del sistema.

\begin{itemize}
  \item \textbf{Automatización Robótica de Procesos (RPA)}
    \begin{itemize}
      \item Definición, evolución y componentes de un sistema RPA.
      \item Criterios de selección de procesos candidatos a automatización.
      \item Ciclo de vida de un proyecto de automatización.
      \item Beneficios, retos y tendencias de investigación en RPA.
    \end{itemize}

  \item \textbf{Arquitecturas de automatización}
    \begin{itemize}
      \item Arquitecturas híbridas: combinación de flujos en la nube y flujos de escritorio.
      \item Automatización de interfaces gráficas de sistemas empresariales (enfoques pro-code, low-code y no-code).
      \item Automatización de hojas de cálculo mediante lenguajes de scripting.
    \end{itemize}

  \item \textbf{Inteligencia de negocios y analítica de datos}
    \begin{itemize}
      \item Definición y alcance de la inteligencia de negocios en la transformación digital.
      \item Proceso analítico: transformación, modelado y cálculo de datos.
      \item Visualización de datos y diseño de dashboards para el monitoreo operativo.
      \item Navegación dinámica en reportes analíticos mediante marcadores y navegadores de botones.
    \end{itemize}

  \item \textbf{Gestión documental}
    \begin{itemize}
      \item Ciclo de vida de un documento en repositorios digitales.
      \item Organización documental: tipos de contenido, metadatos, sitios y bibliotecas.
      \item Automatización de flujos documentales.
    \end{itemize}

  \item \textbf{Sistemas empresariales e integración}
    \begin{itemize}
      \item Definición, módulos y arquitectura de un sistema ERP.
      \item Métodos de integración entre sistemas empresariales (punto a punto, bus de servicios, iPaaS).
      \item Conciliación e interoperabilidad de datos entre plataformas con formatos distintos.
    \end{itemize}

  \item \textbf{Metodología de desarrollo}
    \begin{itemize}
      \item Fases de un proyecto de automatización: planeación, diseño, construcción, prueba y despliegue.
      \item Documentación de bitácoras y coordinación de pruebas.
      \item Control de versiones y respaldo de entregables.
    \end{itemize}
\end{itemize}

\begin{guia}[Guía: Marco teórico (Revisión 2)]
\textbf{¿Para qué sirve?} Reúne los conceptos, teorías, tecnologías, métodos y normas que se necesitan para entender y
justificar tu propuesta. No es un glosario ni un resumen de Internet: cada tema se explica \emph{en función de tu proyecto}.

\textbf{Debe responder (por cada tema):}
\begin{itemize}[nosep,leftmargin=*]
  \item ¿Qué es y cómo funciona? (definición con fuente citada y, si ayuda, un esquema o ecuación).
  \item ¿Para qué se usa en mi proyecto y por qué lo elegí frente a las alternativas?
  \item ¿Qué ventajas, limitaciones o buenas prácticas debo conocer?
\end{itemize}
Además: ¿qué metodología de desarrollo o investigación usarás y por qué? ¿qué normas o estándares aplican (calidad,
seguridad, protección de datos)?

\textbf{Organización:} una \verb|\subsection| por tema, de lo general a lo particular.

\textbf{Extensión mínima:} 400 palabras de desarrollo (sin contar los temas tentativos de la Revisión 1).

\textbf{Fuentes:} libros, artículos y documentación oficial. Cada afirmación importante lleva \verb|\cite{clave}|.

\textbf{Errores frecuentes:} temas sin relación con el proyecto; copiar y pegar; fuentes que son solo URLs de blogs; no
citar; temas que luego no se usan en el desarrollo.

\textbf{Cómo se revisa:} pertinencia, profundidad suficiente, calidad de las fuentes y que cada tema se vincule con la
propuesta.
\end{guia}

% >>> REVISIÓN 2: ESCRIBE AQUÍ el marco teórico (introducción breve y una \subsection por tema) <<<


\subsection{Temas tentativos del marco teórico} \label{sec:temas_tentativos} % audit:req=temas_marco_teorico

\begin{guia}[Guía: Temas tentativos del marco teórico (Revisión 1)]
\textbf{¿Para qué sirve?} En la primera revisión no se pide el marco teórico completo, sino la lista de los temas que
vas a desarrollar. Así el director puede confirmar que vas bien encaminado.

\textbf{Debe responder:} ¿cuáles son los temas (de 5 a 10) que necesito para sustentar el proyecto? Para cada uno, en una
o dos oraciones: ¿qué es y por qué lo necesito?

\textbf{Formato:} una lista (\verb|\begin{enumerate}| o \verb|itemize|), un tema por elemento.

\textbf{Extensión mínima:} 25 palabras. En la Revisión 2 conviertes cada tema en una subsección desarrollada.

\textbf{Cómo se revisa:} que los temas cubran lo necesario para el proyecto y no incluyan relleno.
\end{guia}

% >>> REVISIÓN 1: ESCRIBE AQUÍ la lista de temas tentativos <<<
